\documentclass[11pt,a4paper]{book}
\usepackage[utf8]{inputenc}
\usepackage[T1]{fontenc}
\usepackage[italian]{babel}
\usepackage{amsmath,amssymb,amsthm}

% Ambiente Esercizio: \begin{esercizio}{ID-IMMUTABILE}{DIFFICOLTA_1_10}
\newenvironment{esercizio}[2]{%
  \par\medskip\noindent\textbf{[#1]} \textit{(Difficoltà base: #2/10)}\par\noindent\ignorespaces
}{\par\medskip}

% Ambiente Soluzione: compilato nel tuo PDF ma ignorato dallo script Python per non fare spoiler nell'app
\newenvironment{soluzione}{%
  \par\noindent\textit{Soluzione (trascritta da quaderno A5):}\par\noindent\ignorespaces
}{\par\bigskip}

\begin{document}

% ==============================================================================
\part{Analisi}
% ==============================================================================

\chapter{Pre-Analisi}

\section{Disequazioni Irrazionali}

\begin{esercizio}{PRE-AN-DISEQ-001}{1.5}
Risolvere in $\mathbb{R}$ la disequazione elementare:
$$\sqrt{x + 3} > 2$$
\end{esercizio}
\begin{soluzione}
% Da compilare dopo aver svolto l'esercizio su carta
\end{soluzione}

\begin{esercizio}{PRE-AN-DISEQ-002}{2.0}
Risolvere in $\mathbb{R}$ la disequazione:
$$\sqrt{2x - 1} \le x - 2$$
\end{esercizio}
\begin{soluzione}
\end{soluzione}

\begin{esercizio}{PRE-AN-DISEQ-003}{3.0}
Risolvere in $\mathbb{R}$ la disequazione irrazionale:
$$\sqrt{2x^2 - 5x + 2} > x - 1$$
\end{esercizio}
\begin{soluzione}
\end{soluzione}

\section{Goniometria}

\begin{esercizio}{PRE-AN-GONIOM-001}{2.0}
Risolvere in $[0, 2\pi)$ l'equazione:
$$2\sin^2(x) - \sin(x) - 1 = 0$$
\end{esercizio}
\begin{soluzione}
\end{soluzione}

\chapter{Analisi 1}

\section{Limiti e Asintotici}

\begin{esercizio}{AN1-LIM-001}{3.0}
Calcolare il seguente limite notevole:
$$\lim_{x \to 0} \frac{\sin(3x) - 3x\cos(x)}{x^3}$$
\end{esercizio}
\begin{soluzione}
\end{soluzione}

% ==============================================================================
\part{Geometria}
% ==============================================================================

\chapter{Pre-Geometria}

\section{Coniche nel Piano}

\begin{esercizio}{PRE-GEO-ANALIT-001}{2.5}
Determinare centro, raggio ed eventuali tangenti condotte dall'origine alla circonferenza:
$$x^2 + y^2 - 6x - 4y + 9 = 0$$
\end{esercizio}
\begin{soluzione}
\end{soluzione}

\chapter{Geometria 1}

\section{Autovalori e Diagonalizzazione}

\begin{esercizio}{GEO1-AUTOVAL-001}{3.5}
Determinare il polinomio caratteristico e gli autovalori della matrice:
$$A = \begin{pmatrix} 2 & 1 & 0 \\ 1 & 2 & 0 \\ 0 & 0 & 3 \end{pmatrix}$$
\end{esercizio}
\begin{soluzione}
\end{soluzione}

% ==============================================================================
\part{Algebra}
% ==============================================================================

\chapter{Algebra}

\section{Gruppo Simmetrico e Permutazioni}

\begin{esercizio}{ALG1-SYM-001}{3.0}
In $S_9$, decomporre in cicli disgiunti, calcolare il periodo e la parità della permutazione:
$$\sigma = (1\,4\,7\,2)(2\,5\,8)(3\,6\,4\,9)$$
\end{esercizio}
\begin{soluzione}
\end{soluzione}

\section{Anelli e Congruenze}

\begin{esercizio}{ALG1-ANELLI-001}{4.5}
Sia $A$ un anello commutativo unitario e $I \subseteq A$ un ideale. Dimostrare formalmente che la relazione $a \equiv b \pmod{I} \iff a - b \in I$ è una congruenza compatibile con somma e prodotto.
\end{esercizio}
\begin{soluzione}
\end{soluzione}

\end{document}
